\documentclass{article}
\usepackage{amsmath, amssymb, amsthm}
\usepackage{booktabs}
\usepackage{geometry}
\usepackage{hyperref}

\geometry{margin=1in}
\hypersetup{colorlinks=true, urlcolor=blue}

% Shared macros for the multi-mode mixed-criticality scheduling papers.
% % Shared macros for the multi-mode mixed-criticality scheduling papers.
% % Shared macros for the multi-mode mixed-criticality scheduling papers.
% \input{macros} from any document in this directory rather than redefining
% these locally -- a macro defined in one paper's preamble and silently
% missing from a sibling's is exactly the class of error that put \round and
% \E undefined in evaluation_methodology.tex.

\newcommand{\N}{\mathbb{N}}
\newcommand{\R}{\mathbb{R}}
\newcommand{\Z}{\mathbb{Z}}
\newcommand{\E}{\mathbb{E}}
\newcommand{\Var}{\mathrm{Var}}
\newcommand{\Prob}{\mathbb{P}}
\newcommand{\Ind}{\mathbf{1}}
\newcommand{\ceil}[1]{\lceil #1 \rceil}
\newcommand{\floor}[1]{\lfloor #1 \rfloor}
\newcommand{\round}[1]{\lfloor #1 \rceil}
 from any document in this directory rather than redefining
% these locally -- a macro defined in one paper's preamble and silently
% missing from a sibling's is exactly the class of error that put \round and
% \E undefined in evaluation_methodology.tex.

\newcommand{\N}{\mathbb{N}}
\newcommand{\R}{\mathbb{R}}
\newcommand{\Z}{\mathbb{Z}}
\newcommand{\E}{\mathbb{E}}
\newcommand{\Var}{\mathrm{Var}}
\newcommand{\Prob}{\mathbb{P}}
\newcommand{\Ind}{\mathbf{1}}
\newcommand{\ceil}[1]{\lceil #1 \rceil}
\newcommand{\floor}[1]{\lfloor #1 \rfloor}
\newcommand{\round}[1]{\lfloor #1 \rceil}
 from any document in this directory rather than redefining
% these locally -- a macro defined in one paper's preamble and silently
% missing from a sibling's is exactly the class of error that put \round and
% \E undefined in evaluation_methodology.tex.

\newcommand{\N}{\mathbb{N}}
\newcommand{\R}{\mathbb{R}}
\newcommand{\Z}{\mathbb{Z}}
\newcommand{\E}{\mathbb{E}}
\newcommand{\Var}{\mathrm{Var}}
\newcommand{\Prob}{\mathbb{P}}
\newcommand{\Ind}{\mathbf{1}}
\newcommand{\ceil}[1]{\lceil #1 \rceil}
\newcommand{\floor}[1]{\lfloor #1 \rfloor}
\newcommand{\round}[1]{\lfloor #1 \rceil}


\title{Evaluation Methodology for Multi-Mode Mixed-Criticality Scheduling}
\author{}
\date{}

\begin{document}

\maketitle

\section{Task Set Generation}

Following DRS (RTSS 2020)~\cite{1a02ee59d96a4af497808eeb685b574c}, utilisation vectors are generated via the Dirichlet-Rescale algorithm.

\subsection{DRS Algorithm}

Given $n$ tasks and target utilisation $U$, the DRS algorithm produces $\mathbf{u} = (u_1, \dots, u_n)$ such that:
\[
\sum_{i=1}^n u_i = U, \quad 0 \leq u_i \leq 1
\]

The algorithm operates via:
\begin{enumerate}
    \item Draw $\mathbf{x} \sim \mathrm{Dirichlet}(1, \dots, 1)$ (uniform on simplex)
    \item Scale: $u_i = U \cdot x_i$
    \item Apply constraints via rescaling to standard simplex
\end{enumerate}

\subsection{Task Parameter Generation}

For each task $\tau_i$, parameters are generated as follows:

\begin{align*}
T_i &\sim \mathrm{LogUniform}(\log(T_{\min}), \log(T_{\max})) \\
D_i &= T_i \\
C_i^{(0)} &= \round(u_i \cdot T_i) \\
\mathrm{BCET}_i &\sim \mathrm{Uniform}(BCET_{\mathrm{min}} \cdot C_i^{(0)}, C_i^{(0)}) \\
C_i^{(k-1)} &= \min(\round(CF \cdot C_i^{(0)}), T_i)
\end{align*}

where:
\begin{itemize}
    \item $T_{\min}$ and $T_{\max}$ are the minimum and maximum period bounds
    \item $BCET_{\mathrm{min}}$ is the minimum BCET fraction (typically $0.8$)
    \item $CF$ is the criticality factor relating normal to degraded budgets
\end{itemize}

\section{Metrics}

\subsection{Degradation Metrics}

\begin{align*}
\textbf{NiD} &\quad \text{Number of times degraded mode is entered} \\
\textbf{TiD} &\quad \text{Fraction of time spent in degraded mode} \\
\textbf{JNE} &\quad \text{Jobs Not Executed (LO jobs dropped)} \\
\textbf{LDM} &\quad \text{Late Degraded Mode (LO jobs that miss deadline in degraded mode)} \\
\textbf{HDM} &\quad \text{High Degradation Mode (HI jobs missing deadline)}
\end{align*}

\subsection{Objective Function}

For $k$-level scheduling, we define the objective function:

\[
\Phi = \alpha \cdot \E[\mathrm{JNE}] + \beta \cdot \E[\mathrm{TiD}] + \gamma \cdot \E[\mathrm{WastedCPU}] + \delta \cdot \E[\mathrm{NiD}]
\]

where:
\begin{itemize}
    \item $\mathrm{JNE}$: Jobs Not Executed (LO jobs dropped)
    \item $\mathrm{TiD}$: Fraction of time spent in degraded mode
    \item $\mathrm{WastedCPU}$: CPU cycles spent on abandoned LO work
    \item $\mathrm{NiD}$: Number of times degraded mode is entered (penalizes frequent mode changes)
\end{itemize}
The coefficients $\alpha, \beta, \gamma, \delta$ are weighting parameters that can be tuned based on
application requirements.

\section{Simulation Parameters}

Common parameter values used in evaluation:

\begin{table}[h]
\centering
\caption{Default Simulation Parameters}
\begin{tabular}{lll}
\toprule
Parameter & Default Value & Description \\
\midrule
$T_{\min}$ & 100 & Minimum task period (ticks) \\
$T_{\max}$ & 10000 & Maximum task period (ticks) \\
$BCET_{\mathrm{min}}$ & 0.8 & Minimum BCET as fraction of $C_i^{(0)}$ \\
$CF$ & See Section~\ref{sec:parameters} & Criticality factor \\
$k$ & 2 & Number of degradation levels \\
\bottomrule
\end{tabular}
\end{table}

\section{Notation Summary}

\begin{table}[h]
\centering
\caption{Notation Reference}
\begin{tabular}{ll}
\toprule
Symbol & Meaning \\
\midrule
$n$ & Number of tasks \\
$U$ & Target utilisation \\
$\mathbf{u}$ & Utilisation vector \\
$T_i$ & Period of task $i$ \\
$D_i$ & Deadline of task $i$ \\
$C_i^{(0)}$ & Normal execution budget \\
$C_i^{(k-1)}$ & HI execution budget \\
$\mathrm{BCET}_i$ & Best-case execution time \\
$U_i^{\mathrm{lo}}, U_i^{\mathrm{hi}}$ & Utilisation in LO/HI modes \\
$FP = 1/N$ & HI-behaviour probability per job \\
$CF$ & Criticality factor \\
$k$ & Number of degradation levels \\
$\tau^{\mathrm{HI}}, \tau^{\mathrm{LO}}$ & HI/LO criticality subsets \\
\bottomrule
\end{tabular}
\end{table}

\section{References}

\bibliographystyle{plain}
\bibliography{project_bib,references}

\end{document}
